\section{Giới thiệu và câu hỏi nghiên cứu}
\label{sec:introduction_and_rq}

\subsection{Bối cảnh và Ý nghĩa Thực tiễn của Bài toán STLF}
Dự báo phụ tải điện ngắn hạn (\textit{Short-Term Load Forecasting} -- STLF), với khoảng thời gian dự báo từ 1 giờ đến 24 giờ tới, là bài toán trung tâm trong công tác quản lý và vận hành hệ thống năng lượng hiện đại~\cite{hong2016probabilistic}. Trong một lưới điện xoay chiều, việc duy trì cân bằng công suất tức thời giữa tổng nguồn phát và tổng phụ tải tiêu thụ:
\begin{equation}
    \sum_{i} P_{\text{gen}, i}(t) = \sum_{j} P_{\text{load}, j}(t) + P_{\text{loss}}(t)
\end{equation}
là điều kiện tiên quyết để giữ vững tần số danh định (50/60~Hz) và bảo đảm tính ổn định động học của toàn hệ sinh thái lưới điện. Do năng lượng điện quy mô lớn đòi hỏi chi phí lưu trữ tích năng cực kỳ tốn kém, mọi sai số trong dự báo phụ tải đều kéo theo tổn thất kinh tế khổng lồ: nếu dự báo thiếu, đơn vị vận hành hệ thống điện (\textit{Independent System Operator} -- ISO) buộc phải khởi động khẩn cấp các tổ máy phát chạy dầu/khí dự phòng (\textit{peaking units}) với giá thị trường giao ngay (\textit{spot price}) cao gấp nhiều lần; ngược lại, nếu dự báo thừa, các nhà máy điện cơ bản buộc phải giảm tải non công suất, gây lãng phí nhiên liệu và phát thải khí nhà kính không đáng có. Nghiên cứu kinh điển của Hong và Fan~\cite{hong2016probabilistic} chỉ ra rằng việc giảm chỉ $1\%$ sai số dự báo phụ tải trong một lưới điện cấp tiểu bang có thể tiết kiệm hàng chục triệu USD chi phí vận hành mỗi năm.

Trong bối cảnh hiện nay, quá trình chuyển dịch năng lượng và biến đổi khí hậu toàn cầu đang đặt bài toán STLF trước những thách thức chưa từng có. Tỷ lệ thâm nhập cao của năng lượng tái tạo phân tán (đặc biệt là điện mặt trời áp mái) tạo nên hiện tượng biến dạng phụ tải ròng (\textit{duck curve}), trong khi sự gia tăng của các đợt sóng nhiệt lịch sử kéo dài đẩy nhu cầu sử dụng các thiết bị điều hòa không khí (\textit{HVAC}) vượt xa các đỉnh phụ tải thông thường.

\subsection{Vấn đề Tồn tại và Khoảng trống Nghiên cứu}
Tại cuộc thi \textit{IISE PG\&E Energy Analytics Challenge 2025}~\cite{pge2025zenodo}, công trình của Roy, Pyltsov và Hu~\cite{roy2025hourly} (arXiv:2505.11390) đã chứng minh một phát hiện thực nghiệm then chốt: các mô hình hồi quy dạng bảng phân đoạn theo giờ (24 mô hình XGBoost độc lập cho 24 giờ trong ngày) hoàn toàn đánh bại các kiến trúc \textit{Deep Learning Transformers} phức tạp (PatchTST~\cite{nie2023time}, Informer~\cite{zhou2021informer}, Autoformer~\cite{wu2021autoformer}). Bài báo gốc đạt RMSE $170.84$~MW và MAPE $5.38\%$ trên tập kiểm tra độc lập Year 3 nhờ ứng dụng PCA trên 5 trạm quan trắc nhiệt độ và 5 trạm bức xạ mặt trời kết hợp biến trễ động học 1 giờ ($\text{Lag1}$).

Tuy nhiên, qua quá trình tái lập độc lập và phân tích sâu sắc, chúng tôi nhận diện ba khoảng trống nghiên cứu cốt lõi chưa được giải quyết trong công trình của Roy et al.~\cite{roy2025hourly}:
\begin{enumerate}
    \item \textbf{Thiếu vắng tri thức nhiệt động lực học công trình (\textit{Absence of Thermodynamic Inductive Bias}):} Phụ tải điện chịu sự chi phối nặng nề nhất bởi hành vi điều hòa không khí của con người. Quan hệ giữa nhiệt độ môi trường và công suất điện tiêu thụ tuân theo đường cong chữ U (\textit{U-shaped curve}) với ngưỡng vật lý cơ sở tiện nghi khoảng $18^\circ\text{C}$. Việc chỉ dùng các thành phần chính PCA tuyến tính thuần túy đã làm phẳng các ngưỡng phi tuyến nhạy cảm này, khiến mô hình gặp khó khăn khi ngoại suy các đỉnh phụ tải cực đoan trong các đợt nắng nóng mùa hè.
    \item \textbf{Sự ngắt quãng giả tạo của các biến lịch rời rạc (\textit{Calendar Discontinuity}):} Sử dụng 12 biến giả nhị phân rời rạc (\textit{one-hot dummy variables}) cho 12 tháng tạo ra các bước nhảy bậc thang gián đoạn tại thời điểm giao thừa giữa ngày 31 tháng 12 và ngày 01 tháng 01, hoàn toàn mâu thuẫn với quy luật thiên văn học trơn tru của quỹ đạo Trái Đất.
    \item \textbf{Bỏ ngỏ tính bù trừ đa dạng sai số hai chế độ vận hành (\textit{Lack of Regime-Aware Model Synergy}):} Bài báo gốc chỉ sử dụng duy nhất một họ mô hình cây XGBoost cho mọi thời điểm. Trên thực tế, mô hình tuyến tính có chính quy hóa (ElasticNet) có tính khái quát hóa cực kỳ ổn định trên các vùng phụ tải nền ban đêm ít biến động; ngược lại, mô hình cây quyết định (XGBoost) lại vượt trội trong việc phân tách các quan hệ phi tuyến vào các giờ cao điểm nắng nóng. Việc kết hợp thông minh hai trường phái mô hình này theo ngữ cảnh thời gian thực là hướng đi chưa từng được khám phá.
\end{enumerate}

\subsection{Mục tiêu và Câu hỏi Nghiên cứu (Research Questions)}
Để giải quyết triệt để ba khoảng trống học thuật nêu trên, nghiên cứu này đặt ra mục tiêu mở rộng bài toán STLF của PG\&E thành một khung kiến trúc ensemble đa tầng nhận biết ngữ cảnh hoàn chỉnh. Nghiên cứu tập trung trả lời ba câu hỏi nghiên cứu cụ thể:

\begin{itemize}
    \item \textbf{RQ1:} Việc bổ sung các đặc trưng nhiệt động lực học Heating/Cooling Degree Days (HDD/CDD tại ngưỡng cơ sở $18^\circ\text{C}$ chuẩn ASHRAE) và chuỗi điều hòa Fourier liên tục vào mô hình 24-XGBoost có cải thiện được độ chính xác (RMSE, MAPE, $R^2$) trên tập dữ liệu kiểm tra năm thứ ba so với baseline của Roy et al. (2025) hay không?
    \item \textbf{RQ2:} Mức độ độc lập và tương quan phần dư sai số giữa mô hình tuyến tính ElasticNet và mô hình phi tuyến dạng cây XGBoost là bao nhiêu, và mỗi mô hình thể hiện ưu thế vượt trội trong những chế độ vận hành cụ thể nào của lưới điện?
    \item \textbf{RQ3:} Kiến trúc Context-Aware Stacking (CAS) đề xuất, thông qua việc điều hòa meta-learner bằng vector ngữ cảnh khí tượng và thời gian thực, có khả năng khai thác tối đa tính bù trừ của các mô hình cơ sở để tạo nên bước nhảy vọt về hiệu năng so với các mô hình đơn lẻ hay không?
\end{itemize}

Để trả lời ba câu hỏi trên, nghiên cứu thiết lập ba giả thuyết khoa học tương ứng: \textbf{H1} (Đặc trưng vật lý công trình HDD/CDD và sóng Fourier giúp giảm phương sai mô hình); \textbf{H2} (ElasticNet và XGBoost có tương quan phần dư thấp $r \ll 1.0$ và phân tầng ưu thế theo chế độ phụ tải nền vs sóng nhiệt); và \textbf{H3} (Meta-learner nhận biết ngữ cảnh thời gian thực sẽ học được quy luật định tuyến động vượt trội hơn bất kỳ giải pháp trung bình tĩnh nào).

\subsection{Tóm tắt Đóng góp Cốt lõi của Bài báo}
Các đóng góp học thuật và thực tiễn chính của nghiên cứu bao gồm:
\begin{enumerate}
    \item \textbf{Tái lập độc lập chuẩn xác:} Tái lập hoàn chỉnh toàn bộ quy trình, cấu hình và kết quả của bài báo gốc (Roy et al., 2025) trên ba giao thức kiểm chứng chéo và tập test năm thứ ba, thiết lập nền tảng đối chuẩn vững chắc.
    \item \textbf{Kỹ nghệ đặc trưng định hướng vật lý (Physics-Informed Features):} Bổ sung chỉ số HDD và CDD (ngưỡng $18^\circ\text{C}$) cùng sóng Fourier liên tục ($\sin/\cos$), giúp giảm ngay RMSE từ $171.66$~MW xuống $160.65$~MW (giảm $11.01$~MW so với baseline tái lập, và giảm $10.19$~MW so với baseline công bố $170.84$~MW của bài báo gốc) trên tập kiểm tra Year 3.
    \item \textbf{Khám phá và lượng hóa hiện tượng Chế độ Vận hành Kép (Dual-Regime Discovery):} Chứng minh tương quan phần dư giữa ElasticNet và XGBoost chỉ đạt $r = 0.6131$. Xác nhận thực nghiệm: ElasticNet dẫn đầu trong chế độ phụ tải nền ban đêm (00:00--06:00: RMSE $131.33$~MW so với $141.37$~MW), trong khi XGBoost áp đảo trong các đợt sóng nhiệt (CDD $> 3^\circ\text{C}$: RMSE $244.16$~MW so với $313.27$~MW).
    \item \textbf{Đề xuất kiến trúc Context-Aware Stacking (CAS):} Thiết kế meta-learner LightGBM nhận đồng thời dự báo OOF của mô hình cơ sở và vector ngữ cảnh môi trường. Đạt kỷ lục mới với RMSE \textbf{141.81~MW}, MAPE \textbf{4.75\%}, và $R^2$ đạt \textbf{0.9136} trên tập Year 3 (giảm \textbf{$-$29.03~MW}, tương đương \textbf{$-$17.0\% RMSE} so với bài báo gốc; kiểm định Diebold-Mariano đạt $p < 10^{-5}$).
    \item \textbf{Định lượng giá trị kinh tế và giải thích minh bạch:} Quy đổi mức giảm $29.03$~MW tương đương giá trị tiết kiệm từ \textbf{5.0 đến 7.6 triệu USD mỗi năm} chi phí mua công suất dự phòng quay cho PG\&E; đồng thời trực quan hóa sự dịch chuyển tầm quan trọng của các đặc trưng theo giờ qua biểu đồ nhiệt SHAP 24 giờ.
\end{enumerate}
